\subsection{First-Class Citizens and Higher-Order Functions}

A Python function is an object. Since it is an object, it can be stored in a variable, placed inside a list or dictionary, passed as an argument, returned from another function, and inspected at runtime.

This is what people usually mean when they say that functions are \textit{first-class citizens}. The phrase sounds abstract, but the mechanism is concrete:

\begin{quote}
A function object can travel through the program like other objects.
\end{quote}

A higher-order function is simply a function that receives another function, returns another function, or both.

\subsubsection{Functions as In-Memory Objects: Passing, Returning, and Storing Subroutine References inside Variables}

A function definition creates a function object and binds a name to it.

\begin{verbatim}
def say_doh():
    return "D'oh!"

print(f"say_doh:       {say_doh}")
print(f"type(say_doh): {type(say_doh)}")
print(f"say_doh():     {say_doh()}")
\end{verbatim}

Possible output:

\begin{verbatim}
say_doh:       <function say_doh at 0x...>
type(say_doh): <class 'function'>
say_doh():     D'oh!
\end{verbatim}

The name \texttt{say\_doh} refers to a function object. Writing \texttt{say\_doh} uses the object. Writing \texttt{say\_doh()} calls the object.

This distinction is essential.

\begin{verbatim}
def say_doh():
    return "D'oh!"

fn = say_doh

print(f"fn:       {fn}")
print(f"type(fn): {type(fn)}")
print(f"fn():     {fn()}")
print()
print(f"fn is say_doh: {fn is say_doh}")
\end{verbatim}

Possible output:

\begin{verbatim}
fn:       <function say_doh at 0x...>
type(fn): <class 'function'>
fn():     D'oh!

fn is say_doh: True
\end{verbatim}

The assignment:

\begin{verbatim}
fn = say_doh
\end{verbatim}

does not call the function. It copies the function reference into another name. Both names now refer to the same function object.

The function object can be inspected.

\begin{verbatim}
def say_doh():
    """Return one short line."""
    return "D'oh!"

fn = say_doh

print(f"type(fn): {type(fn)}")
print()

selected_names = [
    "__call__",
    "__name__",
    "__doc__",
    "__defaults__",
    "__kwdefaults__",
    "__annotations__",
    "__code__",
]

for name in selected_names:
    print(f"{name:<18} {name in dir(fn)}")

print()
print(f"fn.__name__: {fn.__name__!r}")
print(f"fn.__doc__:  {fn.__doc__!r}")
\end{verbatim}

Possible output:

\begin{verbatim}
type(fn): <class 'function'>

__call__          True
__name__          True
__doc__           True
__defaults__      True
__kwdefaults__    True
__annotations__   True
__code__          True

fn.__name__: 'say_doh'
fn.__doc__:  'Return one short line.'
\end{verbatim}

The object remains the same function object even if several names point to it.

Functions can be stored inside a list.

\begin{verbatim}
def quote_one():
    return "D'oh!"

def quote_two():
    return "No soup for you!"

def quote_three():
    return "Nobody expects the Spanish Inquisition!"

actions = [quote_one, quote_two, quote_three]

print(f"actions:       {actions}")
print(f"type(actions): {type(actions)}")
print()

for fn in actions:
    print(f"{fn.__name__:<12} -> {fn()}")
\end{verbatim}

Possible output:

\begin{verbatim}
actions:       [<function quote_one at 0x...>, <function quote_two at 0x...>, <function quote_three at 0x...>]
type(actions): <class 'list'>

quote_one    -> D'oh!
quote_two    -> No soup for you!
quote_three  -> Nobody expects the Spanish Inquisition!
\end{verbatim}

The list contains references to function objects. The loop binds \texttt{fn} to each function object and then calls it with \texttt{fn()}.

The list itself is an ordinary list.

\begin{verbatim}
def quote_one():
    return "D'oh!"

def quote_two():
    return "No soup for you!"

actions = [quote_one, quote_two]

print(f"type(actions): {type(actions)}")
print()

selected_names = [
    "append",
    "extend",
    "pop",
    "__len__",
    "__getitem__",
    "__setitem__",
]

for name in selected_names:
    print(f"{name:<14} {name in dir(actions)}")

print()
print(f"len(actions): {len(actions)}")
print(f"actions[0]:   {actions[0]}")
print(f"actions[0](): {actions[0]()}")
\end{verbatim}

Possible output:

\begin{verbatim}
type(actions): <class 'list'>

append         True
extend         True
pop            True
__len__        True
__getitem__    True
__setitem__    True

len(actions): 2
actions[0]:   <function quote_one at 0x...>
actions[0](): D'oh!
\end{verbatim}

There is no special list of functions type. It is a normal list whose elements happen to be function objects.

Functions can also be stored inside dictionaries.

\begin{verbatim}
def small_quote():
    return "Ni!"

def big_quote():
    return "Chuck Norris can divide by zero."

def answer_quote():
    return "The answer is 42."

commands = {
    "small": small_quote,
    "big": big_quote,
    "answer": answer_quote,
}

for key in ["small", "answer", "big"]:
    fn = commands[key]
    print(f"{key:<8} -> {fn()}")
\end{verbatim}

Possible output:

\begin{verbatim}
small    -> Ni!
answer   -> The answer is 42.
big      -> Chuck Norris can divide by zero.
\end{verbatim}

The dictionary maps strings to function objects. Looking up a key returns a function object, and the returned function object can then be called.

The dictionary is ordinary too.

\begin{verbatim}
def small_quote():
    return "Ni!"

commands = {
    "small": small_quote,
}

print(f"type(commands): {type(commands)}")
print()

selected_names = [
    "keys",
    "values",
    "items",
    "get",
    "__getitem__",
    "__setitem__",
]

for name in selected_names:
    print(f"{name:<14} {name in dir(commands)}")

print()
fn = commands.get("small")

print(f"fn:       {fn}")
print(f"type(fn): {type(fn)}")
print(f"fn():     {fn()}")
\end{verbatim}

Possible output:

\begin{verbatim}
type(commands): <class 'dict'>

keys           True
values         True
items          True
get            True
__getitem__    True
__setitem__    True

fn:       <function small_quote at 0x...>
type(fn): <class 'function'>
fn():     Ni!
\end{verbatim}

A function can also return another function object.

\begin{verbatim}
def quote_one():
    return "D'oh!"

def quote_two():
    return "No soup for you!"

def choose_quote(kind):
    if kind == "simpsons":
        return quote_one

    return quote_two

selected_fn = choose_quote("simpsons")

print(f"selected_fn:       {selected_fn}")
print(f"type(selected_fn): {type(selected_fn)}")
print(f"selected_fn():     {selected_fn()}")
\end{verbatim}

Possible output:

\begin{verbatim}
selected_fn:       <function quote_one at 0x...>
type(selected_fn): <class 'function'>
selected_fn():     D'oh!
\end{verbatim}

The function \texttt{choose\_quote()} does not return the result of calling \texttt{quote\_one}. It returns the function object \texttt{quote\_one} itself.

The difference is visible here:

\begin{verbatim}
def quote_one():
    return "D'oh!"

def return_function_object():
    return quote_one

def return_function_result():
    return quote_one()

a = return_function_object()
b = return_function_result()

print(f"a:       {a}")
print(f"type(a): {type(a)}")
print(f"a():     {a()}")
print()

print(f"b:       {b!r}")
print(f"type(b): {type(b)}")
\end{verbatim}

Possible output:

\begin{verbatim}
a:       <function quote_one at 0x...>
type(a): <class 'function'>
a():     D'oh!

b:       "D'oh!"
type(b): <class 'str'>
\end{verbatim}

The expression \texttt{quote\_one} produces the function object. The expression \texttt{quote\_one()} calls the function and produces its return value.

The first-class rule is therefore not philosophical. It is ordinary object-reference behavior:

\begin{verbatim}
name = function_object
list_item = function_object
dict_value = function_object
return function_object
call function_object later
\end{verbatim}

\subsubsection{Callbacks and Callback Chains: Decoupling Structural Execution Graphs}

A callback is a function object passed into another function so the receiving function can call it later.

The callback function does not control the outer structure. It supplies a piece of behavior that the outer function uses at a chosen point.

\begin{verbatim}
def shout(text):
    return text.upper()

def whisper(text):
    return text.lower()

def process_line(text, callback):
    print(f"original:  {text!r}")

    result = callback(text)

    print(f"processed: {result!r}")
    return result

process_line("No soup for you!", shout)
print()
process_line("No soup for you!", whisper)
\end{verbatim}

Possible output:

\begin{verbatim}
original:  'No soup for you!'
processed: 'NO SOUP FOR YOU!'

original:  'No soup for you!'
processed: 'no soup for you!'
\end{verbatim}

The function \texttt{process\_line()} controls the execution frame:

\begin{verbatim}
print original text
call callback
print processed text
return processed text
\end{verbatim}

The callback controls only the transformation step.

The callback is a normal function object.

\begin{verbatim}
def shout(text):
    return text.upper()

def process_line(text, callback):
    print(f"callback:       {callback}")
    print(f"type(callback): {type(callback)}")
    print(f"callback name:  {callback.__name__}")

    return callback(text)

result = process_line("D'oh!", shout)

print()
print(f"result: {result}")
\end{verbatim}

Possible output:

\begin{verbatim}
callback:       <function shout at 0x...>
type(callback): <class 'function'>
callback name:  shout

result: D'OH!
\end{verbatim}

The parameter name \texttt{callback} is just a local name bound to a function object supplied by the caller.

Callbacks are useful because they separate the stable structure from the replaceable behavior.

\begin{verbatim}
def format_plain(text):
    return text

def format_boxed(text):
    return f"[ {text} ]"

def format_numbered(text):
    return f"42. {text}"

def print_report_line(text, formatter):
    formatted = formatter(text)
    print(formatted)

print_report_line("It is only a model.", format_plain)
print_report_line("It is only a model.", format_boxed)
print_report_line("It is only a model.", format_numbered)
\end{verbatim}

Possible output:

\begin{verbatim}
It is only a model.
[ It is only a model. ]
42. It is only a model.
\end{verbatim}

The outer function does not need \texttt{if} branches for every formatting strategy. It receives the strategy as a function object.

A callback chain applies several callback functions in order.

\begin{verbatim}
def strip_text(text):
    return text.strip()

def upper_text(text):
    return text.upper()

def add_suffix(text):
    return text + "!"

def run_chain(text, callbacks):
    current = text

    for callback in callbacks:
        print(f"before {callback.__name__:<12}: {current!r}")
        current = callback(current)

    print(f"final result:       {current!r}")
    return current

steps = [strip_text, upper_text, add_suffix]

run_chain("  ni  ", steps)
\end{verbatim}

Possible output:

\begin{verbatim}
before strip_text  : '  ni  '
before upper_text  : 'ni'
before add_suffix  : 'NI'
final result:       'NI!'
\end{verbatim}

The list \texttt{steps} stores function objects. The function \texttt{run\_chain()} walks through that list and calls each function object in sequence.

This builds a small execution graph:

\begin{verbatim}
input text
    -> strip_text
    -> upper_text
    -> add_suffix
    -> final output
\end{verbatim}

The chain itself is data. It can be changed without rewriting \texttt{run\_chain()}.

\begin{verbatim}
def strip_text(text):
    return text.strip()

def upper_text(text):
    return text.upper()

def lower_text(text):
    return text.lower()

def add_suffix(text):
    return text + "!"

def run_chain(text, callbacks):
    current = text

    for callback in callbacks:
        current = callback(current)

    return current

loud_chain = [strip_text, upper_text, add_suffix]
quiet_chain = [strip_text, lower_text]

print(run_chain("  No Soup For You  ", loud_chain))
print(run_chain("  No Soup For You  ", quiet_chain))
\end{verbatim}

Possible output:

\begin{verbatim}
NO SOUP FOR YOU!
no soup for you
\end{verbatim}

The code that runs the chain is unchanged. Only the list of function objects changes.

The chain list can be inspected.

\begin{verbatim}
def strip_text(text):
    return text.strip()

def upper_text(text):
    return text.upper()

def add_suffix(text):
    return text + "!"

steps = [strip_text, upper_text, add_suffix]

print(f"type(steps): {type(steps)}")
print()

for index, fn in enumerate(steps):
    print(f"index={index}")
    print(f"  object: {fn}")
    print(f"  type:   {type(fn)}")
    print(f"  name:   {fn.__name__}")
\end{verbatim}

Possible output:

\begin{verbatim}
type(steps): <class 'list'>

index=0
  object: <function strip_text at 0x...>
  type:   <class 'function'>
  name:   strip_text
index=1
  object: <function upper_text at 0x...>
  type:   <class 'function'>
  name:   upper_text
index=2
  object: <function add_suffix at 0x...>
  type:   <class 'function'>
  name:   add_suffix
\end{verbatim}

Callbacks can also be selected from a dictionary.

\begin{verbatim}
def as_upper(text):
    return text.upper()

def as_lower(text):
    return text.lower()

def as_title(text):
    return text.title()

formatters = {
    "upper": as_upper,
    "lower": as_lower,
    "title": as_title,
}

mode = "title"
text = "these pretzels are making me thirsty"

formatter = formatters[mode]
result = formatter(text)

print(f"mode:      {mode}")
print(f"formatter: {formatter.__name__}")
print(f"result:    {result}")
\end{verbatim}

Possible output:

\begin{verbatim}
mode:      title
formatter: as_title
result:    These Pretzels Are Making Me Thirsty
\end{verbatim}

This dictionary is a dispatch table. It replaces a chain of conditionals with a mapping from command names to function objects.

A simple conditional version would be:

\begin{verbatim}
if mode == "upper":
    result = as_upper(text)
elif mode == "lower":
    result = as_lower(text)
elif mode == "title":
    result = as_title(text)
\end{verbatim}

The dispatch-table version is:

\begin{verbatim}
formatter = formatters[mode]
result = formatter(text)
\end{verbatim}

Both are valid. The dispatch table makes the executable choices explicit as data.

A safer version checks whether the requested key exists.

\begin{verbatim}
def as_upper(text):
    return text.upper()

def as_lower(text):
    return text.lower()

formatters = {
    "upper": as_upper,
    "lower": as_lower,
}

mode = "reverse"
text = "D'oh!"

formatter = formatters.get(mode)

if formatter is None:
    print(f"unknown mode: {mode!r}")
else:
    print(formatter(text))
\end{verbatim}

Possible output:

\begin{verbatim}
unknown mode: 'reverse'
\end{verbatim}

The value returned by \texttt{dict.get()} is either a function object or \texttt{None}. The call is performed only after checking that a function object was actually found.

The core callback rule is:

\begin{quote}
A callback is a function object supplied as data so another function can call it at a controlled point in its own execution.
\end{quote}

\subsubsection{Higher-Order Functions: Functions That Receive or Return Other Functions}

A higher-order function is a function that works with other functions as values.

It may receive a function:

\begin{verbatim}
def double(x):
    return x * 2

def apply_once(fn, value):
    return fn(value)

result = apply_once(double, 21)

print(f"result: {result}")
\end{verbatim}

Possible output:

\begin{verbatim}
result: 42
\end{verbatim}

The function \texttt{apply\_once()} receives the function object \texttt{double}, binds it to its local name \texttt{fn}, and calls it.

The binding is ordinary parameter binding.

\begin{verbatim}
def double(x):
    return x * 2

def apply_once(fn, value):
    print(f"fn:       {fn}")
    print(f"type(fn): {type(fn)}")
    print(f"value:    {value}")
    print()

    return fn(value)

result = apply_once(double, 21)

print(f"result: {result}")
\end{verbatim}

Possible output:

\begin{verbatim}
fn:       <function double at 0x...>
type(fn): <class 'function'>
value:    21

result: 42
\end{verbatim}

A higher-order function may apply a function to several values.

\begin{verbatim}
def square(x):
    return x * x

def apply_to_all(fn, values):
    results = []

    for value in values:
        result = fn(value)
        results.append(result)

    return results

nums = [1, 2, 3, 4]

squares = apply_to_all(square, nums)

print(f"nums:    {nums}")
print(f"squares: {squares}")
\end{verbatim}

Possible output:

\begin{verbatim}
nums:    [1, 2, 3, 4]
squares: [1, 4, 9, 16]
\end{verbatim}

The result list is a normal list.

\begin{verbatim}
def square(x):
    return x * x

def apply_to_all(fn, values):
    results = []

    for value in values:
        results.append(fn(value))

    return results

squares = apply_to_all(square, [1, 2, 3, 4])

print(f"type(squares): {type(squares)}")
print()

selected_names = [
    "append",
    "extend",
    "pop",
    "__len__",
    "__getitem__",
]

for name in selected_names:
    print(f"{name:<14} {name in dir(squares)}")

print()
print(f"len(squares): {len(squares)}")
print(f"squares[2]:   {squares[2]}")
\end{verbatim}

Possible output:

\begin{verbatim}
type(squares): <class 'list'>

append         True
extend         True
pop            True
__len__        True
__getitem__    True

len(squares): 4
squares[2]:   9
\end{verbatim}

The higher-order structure is:

\begin{verbatim}
apply_to_all receives:
    fn     -> function object
    values -> iterable object

inside the loop:
    call fn(value)
    collect returned object
\end{verbatim}

A higher-order function may also return a function.

\begin{verbatim}
def make_multiplier(factor):
    def multiply(x):
        return x * factor

    return multiply

times_two = make_multiplier(2)
times_three = make_multiplier(3)

print(f"times_two:   {times_two}")
print(f"times_three: {times_three}")
print()
print(f"times_two(21):   {times_two(21)}")
print(f"times_three(14): {times_three(14)}")
\end{verbatim}

Possible output:

\begin{verbatim}
times_two:   <function make_multiplier.<locals>.multiply at 0x...>
times_three: <function make_multiplier.<locals>.multiply at 0x...>

times_two(21):   42
times_three(14): 42
\end{verbatim}

The function \texttt{make\_multiplier()} returns the inner function object \texttt{multiply}. The returned function can be stored and called later.

This example also involves a closure because the returned function uses \texttt{factor} after the outer function has returned. Closure internals will be studied in the next subsection. For now, the important first-class behavior is that a function can return another function object.

A simpler version can return one of several existing functions without introducing closure state.

\begin{verbatim}
def as_upper(text):
    return text.upper()

def as_lower(text):
    return text.lower()

def choose_formatter(mode):
    if mode == "upper":
        return as_upper

    return as_lower

formatter = choose_formatter("upper")

print(f"formatter:       {formatter}")
print(f"type(formatter): {type(formatter)}")
print(f"result:          {formatter('ni!')}")
\end{verbatim}

Possible output:

\begin{verbatim}
formatter:       <function as_upper at 0x...>
type(formatter): <class 'function'>
result:          NI!
\end{verbatim}

The returned object is selected at runtime. The caller receives a function object and can call it later.

A function can both receive and return functions.

\begin{verbatim}
def shout(text):
    return text.upper()

def make_logged(fn):
    def logged(text):
        print(f"calling {fn.__name__} with {text!r}")
        result = fn(text)
        print(f"result is {result!r}")
        return result

    return logged

logged_shout = make_logged(shout)

final_text = logged_shout("no soup for you!")

print(f"final_text: {final_text}")
\end{verbatim}

Possible output:

\begin{verbatim}
calling shout with 'no soup for you!'
result is 'NO SOUP FOR YOU!'
final_text: NO SOUP FOR YOU!
\end{verbatim}

The function \texttt{make\_logged()} receives one function object and returns another function object. The returned function calls the original function and adds behavior around it.

This is the basic shape behind decorators, but the decorator syntax itself will be explained later. Here, the mechanism is still plain function-object passing and returning.

The returned function object can be inspected.

\begin{verbatim}
def shout(text):
    return text.upper()

def make_logged(fn):
    def logged(text):
        print(f"calling {fn.__name__}")
        return fn(text)

    return logged

logged_shout = make_logged(shout)

print(f"logged_shout:       {logged_shout}")
print(f"type(logged_shout): {type(logged_shout)}")
print(f"name:               {logged_shout.__name__}")
print(f"closure:            {logged_shout.__closure__}")
\end{verbatim}

Possible output:

\begin{verbatim}
logged_shout:       <function make_logged.<locals>.logged at 0x...>
type(logged_shout): <class 'function'>
name:               logged
closure:            (<cell at 0x...: function object at 0x...>,)
\end{verbatim}

The closure appears because \texttt{logged()} remembers \texttt{fn}. The detailed mechanics of \texttt{\_\_closure\_\_} and cell objects belong to the next subsection.

For now, the architectural point is enough:

\begin{quote}
Higher-order functions let a program pass executable behavior through ordinary object references.
\end{quote}

This allows code structure to be separated from the specific operation performed at a given step.

A compact summary:

\begin{center}
\begin{tabular}{|p{0.30\linewidth}|p{0.52\linewidth}|}
\hline
\textbf{Pattern} & \textbf{Meaning} \\
\hline
\texttt{fn = some\_function} & store a function object in another name \\
\hline
\texttt{items = [fn1, fn2]} & store function objects inside a list \\
\hline
\texttt{table["x"] = fn} & store a function object inside a dictionary \\
\hline
\texttt{outer(fn)} & pass a function object into another function \\
\hline
\texttt{return fn} & return a function object from a function \\
\hline
\texttt{fn()} & call the function object currently referenced by \texttt{fn} \\
\hline
\end{tabular}
\end{center}

The final rule is:

\begin{quote}
A function name without parentheses refers to the function object. A function name with parentheses calls the function object and produces its return value.
\end{quote}